Early uncertainty collapse can discard correct transcription alternatives before grouped or downstream constraints can use them. We study this problem in low-resource glyph and grouped recognition with an auditable uncertainty-preserving pipeline built around confusion-network posteriors, constrained decoding, and explicit failure analysis. The evidence spans three synthetic-from-real glyph benchmarks, two real grouped token-aligned corpora, and a redesigned real downstream task with substantially improved support coverage.

Across Omniglot, scikit-learn Digits, and Kuzushiji-49, preserving uncertainty improves symbol retention and yields modest grouped gains, with grouped top-k improvements more stable than exact-match gains. On Historical Newspapers and ScaDS.AI handwriting, grouped top-k rescue and symbol top-k rescue both transfer, establishing a replicated real grouped result. Historical Newspapers also survives a full-split audit and gold-style adjudication with only 2 of 126 token corrections and no metric drift. Exact grouped and downstream recovery remain mixed: the redesigned real downstream task reveals selective gains, including conformal gains on Historical Newspapers and raw uncertainty gains on ScaDS.AI, but not a clean replicated exact win.

The paper's contribution is explanatory rather than architectural. A support-aware propagation analysis shows that symbol rescue is the strongest predictor of grouped rescue, while grouped rescue is necessary but not sufficient for downstream success. Real downstream rescue becomes much more likely once grouped top-k gain reaches at least one recovered alternative, and conformal helps in a different support regime than raw uncertainty. The resulting claim is intentionally narrow: preserved uncertainty reliably improves local and grouped alternative retention, but higher-level usefulness remains support-gated.
